\slajdid{09_1-s041}
\begin{frame}[label=09_1-s041]
\gusttekst
\[k_{n,D}(V_{DD}-V_O)(V_{DD}+2\Delta-V_O)=k_{n,L}(V_{DD}+\Delta-V_O)^2\]
Zanimljivo je što za $\Delta=0$ kada je tranzistor $T_L$ na prelazu između omske oblasti i zasićenja rešenje jednačine je uz $k_{n,L}\ne k_{n,D}$ moguće samo za $V_O=V_{DD}$ pa je tada i $V_{IL}=V_{Tn,D}$ što ćemo videti i kasnije kada budemo analizirali situaciju kada tranzistor $T_L$ uvek radi u zasićenju.\par\medskip
Opet je slična situacija, zgodno je pisati po $V_{DD}-V_O$
\[k_{n,D}(V_{DD}-V_O)^2+2\Delta k_{n,D}(V_{DD}-V_O)=k_{n,L}\bigl((V_{DD}-V_O)^2+2\Delta(V_{DD}-V_O)+\Delta^2\bigr)\]
\[(V_{DD}-V_O)^2(k_{n,D}-k_{n,L})+2\Delta(V_{DD}-V_O)(k_{n,D}-k_{n,L})-\Delta^2 k_{n,L}=0\]
\par\bigskip
Ono što odmah treba uočiti jeste na primer ako je $k_{n,L}=k_{n,D}$ jednačina nema rešenje. A to znači da u ovoj oblasti nema tačke gde je pojačanje jednako jedinici.
\end{frame}
